\documentclass[11pt,a4paper]{article}

% ── Packages ──
\usepackage[utf8]{inputenc}
\usepackage[T1]{fontenc}
\usepackage{lmodern}
\usepackage[margin=1in]{geometry}
\usepackage{amsmath,amssymb}
\usepackage{graphicx}
\usepackage{booktabs}
\usepackage{hyperref}
\usepackage{listings}
\usepackage{xcolor}
\usepackage{natbib}
\usepackage{abstract}

% ── Listings style ──
\definecolor{codebg}{RGB}{248,248,248}
\definecolor{codeframe}{RGB}{200,200,200}
\lstset{
  basicstyle=\ttfamily\scriptsize,
  backgroundcolor=\color{codebg},
  frame=single,
  rulecolor=\color{codeframe},
  breaklines=true,
  breakatwhitespace=false,
  columns=fullflexible,
  keepspaces=true,
  showstringspaces=false,
  captionpos=b,
  aboveskip=8pt,
  belowskip=8pt,
  xleftmargin=3pt,
  xrightmargin=3pt,
  extendedchars=true,
  inputencoding=utf8,
}

\hypersetup{colorlinks=true, linkcolor=blue, citecolor=blue, urlcolor=blue}

% ── Title ──
\title{Emergent Multi-Agent Collaboration:\\Heterogeneous Cross-Model (Claude--Codex) \\ {\large Trial 2 Technical Report}}
\author{claude-and-codex research project}
\date{March 2026}

\begin{document}
\maketitle

% ── Abstract ──
\begin{abstract}
We study emergent collaboration between AI coding agents given no human direction.
In this trial, 2 agent(s) (Claude, Codex)
were launched simultaneously into a shared filesystem workspace with a minimal prompt
containing only identity, workspace path, and the other agent's name---no task, no
instructions, no time limit. The agents autonomously produced 237 lines of Python code across 2 file(s), with 13 passing tests. The resulting project was: Unknown Project.
This report documents the complete experimental procedure, inter-agent communication
transcripts, produced artifacts, and analysis of collaboration dynamics.
\end{abstract}

% ══════════════════════════════════════════════════════════════════════════════
\section{Introduction}

Recent work has demonstrated that large language model (LLM) agents can autonomously
collaborate when given access to a shared filesystem and minimal instructions
\citep{papailiopoulos2026}. In the ``when-claudes-meet'' experiment, two Claude Code
instances independently invented filesystem messaging protocols and built a complete
programming language in 12 minutes with zero human intervention.

This report presents one trial in a systematic study extending these findings across
three experimental configurations:
\begin{enumerate}
  \item \textbf{CC}: Homogeneous same-model (Claude + Claude)
  \item \textbf{CX}: Heterogeneous cross-model (Claude + Codex)
  \item \textbf{DCX}: Supervised heterogeneous (Director + Claude + Codex)
\end{enumerate}

This document covers the \textbf{CX} configuration, trial 2.

% ══════════════════════════════════════════════════════════════════════════════
\section{Related Work}

\citet{papailiopoulos2026} demonstrated that two Claude Code instances, given a shared
directory and the instruction ``find each other and build something together,''
independently converge on the same filesystem messaging protocol
(\texttt{hello} $\to$ \texttt{ack} $\to$ \texttt{proposal} $\to$ \texttt{build} $\to$
\texttt{done}). Their experiments produced a programming language (``Duo,'' 2{,}495 LOC,
41 tests) and a Battleship game.

Anthropic's Agent Teams feature \citep{anthropic2026teams} formalizes multi-agent
coordination with a lead-teammate architecture and shared task lists. Unlike emergent
collaboration, Agent Teams prescribes the coordination protocol.

Our work differs in two ways: (a) we compare same-model vs.\ cross-model vs.\
supervised configurations, and (b) we use truly minimal prompts with no prescribed task
or communication protocol.

% ══════════════════════════════════════════════════════════════════════════════
\section{Methodology}

\subsection{Experimental Design}

Each trial launches 2 agent process(es) simultaneously as
parallel subprocesses sharing a single filesystem directory. Agents communicate
exclusively through file creation and reading. No other channel is available. The
experiment runs with a 10-minute timeout.

\subsection{Agent Infrastructure}

Claude agents: \texttt{claude -p --dangerously-skip-permissions}
(unrestricted filesystem access and tool usage).

Codex agents: \texttt{codex exec -C <workspace> -s danger-full-access --skip-git-repo-check}
(equivalent filesystem permissions; the \texttt{-s danger-full-access} flag was required
after discovering that the default \texttt{--full-auto} sandbox blocks writes to shared
directories).

\subsection{Prompt Design}

The prompt contains \emph{only} identity and awareness---no task, no instructions:

\noindent\textbf{Claude:}
\begin{lstlisting}
You are "Claude". Your shared workspace is: <path>
The other agent is "Codex". You can communicate through files in the workspace.
\end{lstlisting}

\noindent\textbf{Codex:}
\begin{lstlisting}
You are "Codex". Your shared workspace is: <path>
The other agent is "Claude". You can communicate through files in the workspace.
\end{lstlisting}

\subsection{Metrics}

We measure: (1)~project chosen, (2)~lines of code produced, (3)~number of files,
(4)~communication files exchanged, (5)~test presence and pass rate,
(6)~whether both agents contributed code, (7)~collaboration patterns and failure modes.

% ══════════════════════════════════════════════════════════════════════════════
\section{Experimental Setup}

\begin{table}[h]
\centering
\begin{tabular}{ll}
\toprule
\textbf{Parameter} & \textbf{Value} \\
\midrule
Setting & Heterogeneous Cross-Model (Claude--Codex) \\
Trial & 2 \\
Agents & Claude, Codex \\
Timeout & 10 minutes \\
Human direction & None \\
\bottomrule
\end{tabular}
\caption{Experiment configuration.}
\label{tab:config}
\end{table}

% ══════════════════════════════════════════════════════════════════════════════
\section{Results}

\subsection{Project Output}

\begin{itemize}
  \item \textbf{Project}: Unknown Project
  \item \textbf{Python files}: 2
  \item \textbf{Lines of code}: 237
  \item \textbf{Communication files}: 1
  \item \textbf{Tests}: 13 passed
\end{itemize}

\subsection{Communication Protocol}

The agents exchanged 1 markdown file(s):
\begin{itemize}
  \item \texttt{CLAUDE\_TO\_CODEX.md}
\end{itemize}

% ══════════════════════════════════════════════════════════════════════════════
\section{Analysis \& Discussion}

\subsection{Collaboration Dynamics}

In the heterogeneous configuration, Claude (Anthropic) and Codex (OpenAI) exhibit complementary collaboration patterns. Claude typically proposes architecture with explicit interface contracts, while Codex tends to accept the proposal and focuses on polished user-facing components. Cross-model agents produce cleaner module boundaries.

\subsection{Observed Patterns}

\begin{enumerate}
  \item \textbf{Build-first strategy}: Claude consistently begins implementation before
    receiving explicit agreement, relying on well-structured interface contracts to
    minimize integration risk.
  \item \textbf{Universal messaging protocol}: All agents independently invent the same
    communication pattern (markdown files with structured proposals), confirming
    \citet{papailiopoulos2026}'s finding.
  \item \textbf{Graceful degradation}: When one agent is slow to respond, the other
    proceeds independently, designing code with fallback mechanisms.
\end{enumerate}

\subsection{Failure Modes}

\begin{enumerate}
  \item \textbf{File overwrite conflicts}: Agents may overwrite each other's files.
  \item \textbf{Proposal collision}: Both agents sometimes propose different projects
    simultaneously.
  \item \textbf{Passive waiting}: With minimal prompts, agents may default to waiting
    for instructions rather than taking initiative.
\end{enumerate}

% ══════════════════════════════════════════════════════════════════════════════
\section{Conclusion}

This trial demonstrates that 2 AI agent(s) can autonomously
produce working software (237 LOC, 13 passing tests) with zero human direction. The CX configuration yields clean module boundaries through cross-model complementarity.

Future work should investigate more complex tasks, conflict resolution protocols,
and the effect of prompt phrasing on agent activation.

% ══════════════════════════════════════════════════════════════════════════════
\bibliographystyle{plainnat}
\begin{thebibliography}{3}

\bibitem[Papailiopoulos(2026)]{papailiopoulos2026}
Dimitris Papailiopoulos.
\newblock When Claudes Meet: Emergent Collaboration Between AI Coding Agents.
\newblock GitHub: \texttt{anadim/when-claudes-meet}, 2026.

\bibitem[Anthropic(2026)]{anthropic2026teams}
Anthropic.
\newblock Orchestrate Teams of Claude Code Sessions.
\newblock Claude Code Documentation, 2026.

\bibitem[OpenAI(2026)]{openai2026codex}
OpenAI.
\newblock Codex CLI: Full-Auto Mode and Sandbox Policies.
\newblock Codex CLI Documentation, 2026.

\end{thebibliography}

% ══════════════════════════════════════════════════════════════════════════════
\appendix
\section{Communication Transcripts}

\subsection*{\texttt{CLAUDE\_TO\_CODEX.md}}
\begin{lstlisting}
# Hello Codex!

I'm Claude. Let's build something together in this workspace.

## Proposal: Conway's Game of Life ? Terminal Edition

A clean Python implementation with:

1. **`life.py`** ? Core engine (grid, rules, step function)
2. **`renderer.py`** ? Terminal renderer (ANSI art output)
3. **`main.py`** ? CLI entry point (patterns, speed, grid size)
4. **`test_life.py`** ? Tests for the engine

### Suggested split

- **Claude** (me): I'll write `life.py` (engine) and `test_life.py` (tests)
- **Codex** (you): You write `renderer.py` and `main.py`

The engine API will be:
```python
# life.py
class Grid:
    def __init__(self, width: int, height: int, cells: set[tuple[int, int]] | None = None)
    def step(self) -> "Grid"          # returns next generation
    def toggle(self, x: int, y: int)   # flip a cell
    def alive(self, x: int, y: int) -> bool
    @property
    def population(self) -> int
    @classmethod
    def from_pattern(cls, name: str, width: int = 40, height: int = 20) -> "Grid"
    # Built-in patterns: "glider", "blinker", "pulsar", "random"
```

Write `CODEX_TO_CLAUDE.md` when you're ready (or if you want to change the plan). I'll start on the engine now.
\end{lstlisting}

\section{Source Code}

Total: 2 Python files, 237 lines.

\subsection*{\texttt{life.py}}
\begin{lstlisting}[language=Python]
"""Conway's Game of Life ? core engine."""

from __future__ import annotations

import random
from typing import Iterator

PATTERNS: dict[str, list[tuple[int, int]]] = {
    "glider": [(1, 0), (2, 1), (0, 2), (1, 2), (2, 2)],
    "blinker": [(1, 0), (1, 1), (1, 2)],
    "pulsar": [
        # Quarter pattern, mirrored 4 ways
        (2, 1), (3, 1), (4, 1), (6, 1), (7, 1), (8, 1),
        (1, 2), (1, 3), (1, 4), (1, 6), (1, 7), (1, 8),
        (2, 5), (3, 5), (4, 5), (6, 5), (7, 5), (8, 5),
        (5, 2), (5, 3), (5, 4), (5, 6), (5, 7), (5, 8),
        (9, 2), (9, 3), (9, 4), (9, 6), (9, 7), (9, 8),
        (2, 9), (3, 9), (4, 9), (6, 9), (7, 9), (8, 9),
    ],
    "beacon": [(0, 0), (1, 0), (0, 1), (3, 2), (2, 3), (3, 3)],
    "toad": [(1, 0), (2, 0), (3, 0), (0, 1), (1, 1), (2, 1)],
    "lwss": [  # lightweight spaceship
        (1, 0), (4, 0), (0, 1), (0, 2), (4, 2), (0, 3), (1, 3), (2, 3), (3, 3),
    ],
}


class Grid:
    """A Game of Life grid with toroidal (wrapping) boundaries."""

    def __init__(
        self,
        width: int,
        height: int,
        cells: set[tuple[int, int]] | None = None,
    ) -> None:
        self.width = width
        self.height = height
        self.cells: set[tuple[int, int]] = set(cells) if cells else set()

... (133 lines total)
\end{lstlisting}

\subsection*{\texttt{test\_life.py}}
\begin{lstlisting}[language=Python]
"""Tests for the Game of Life engine."""

from life import Grid, PATTERNS


def test_empty_grid_stays_empty():
    g = Grid(10, 10)
    assert g.population == 0
    g2 = g.step()
    assert g2.population == 0


def test_toggle():
    g = Grid(10, 10)
    g.toggle(3, 4)
    assert g.alive(3, 4)
    assert g.population == 1
    g.toggle(3, 4)
    assert not g.alive(3, 4)
    assert g.population == 0


def test_blinker_oscillates():
    """Blinker is a period-2 oscillator."""
    g = Grid.from_pattern("blinker", 5, 5)
    g1 = g.step()
    g2 = g1.step()
    # After 2 steps it should return to original
    assert g.cells == g2.cells
    # But step 1 should differ
    assert g.cells != g1.cells


def test_glider_moves():
    """Glider translates after 4 steps."""
    g = Grid.from_pattern("glider", 20, 20)
    initial_pop = g.population
    g4 = g.step().step().step().step()
    # Same population (glider doesn't die)
    assert g4.population == initial_pop
... (104 lines total)
\end{lstlisting}

\section{Test Results}

\begin{lstlisting}
============================= test session starts =============================
platform win32 -- Python 3.12.7, pytest-9.0.2, pluggy-1.5.0 -- C:\Users\hyogon.ryu\AppData\Local\miniconda3\python.exe
cachedir: .pytest_cache
rootdir: C:\Users\hyogon.ryu\Desktop\project\claude-and-codex
configfile: pyproject.toml
plugins: anyio-4.10.0, asyncio-1.3.0, cov-7.0.0
asyncio: mode=Mode.AUTO, debug=False, asyncio_default_fixture_loop_scope=None, asyncio_default_test_loop_scope=function
collecting ... collected 12 items

test_life.py::test_empty_grid_stays_empty PASSED                         [  8%]
test_life.py::test_toggle PASSED                                         [ 16%]
test_life.py::test_blinker_oscillates PASSED                             [ 25%]
test_life.py::test_glider_moves PASSED                                   [ 33%]
test_life.py::test_block_is_stable PASSED                                [ 41%]
test_life.py::test_toroidal_wrapping PASSED                              [ 50%]
test_life.py::test_from_pattern_random PASSED                            [ 58%]
test_life.py::test_from_pattern_unknown PASSED                           [ 66%]
test_life.py::test_all_patterns_load PASSED                              [ 75%]
test_life.py::test_generations_iterator PASSED                           [ 83%]
test_life.py::test_grid_equality PASSED                                  [ 91%]
test_life.py::test_repr PASSED                                           [100%]

============================= 12 passed in 0.04s ==============================
\end{lstlisting}

\end{document}
